% =============================================================================
% Iter 43 -- Pillar 3: G=4 vs G=32 ZVF-Decomposed + Cost-Adjusted Equivalence
% Driver: scripts/group_size_iter43.py + scripts/group_size_iter43_fig.py
% Inputs: experiments/results/group_size_token_normalized.tsv
% Outputs: experiments/results/group_size_iter43_{eff_zvf,difficulty,flop_tost,summary}.tsv
%          figures/group_size_iter43.{pdf,png}
% =============================================================================
\subsection{Iter 43 Pillar 3: ZVF-Decomposed Equivalence at Broad Scale}
\label{sec:group-size-iter43}

\paragraph{Question.}
Iter 31's iso-token battery showed Wu~2025's 97.6\% retention claim
fails on $G{=}4$ vs $G{=}32$ at every budget $\geq 4$\,M; iter 35's
40-cell pair sweep confirmed a population-level pass rate of 18/40;
iter 39 sharpened this into a critical-budget $T^*{=}0.61$\,M with
bootstrap 95\% CI $[0.32, 3.45]$\,M and showed the iter 35 pass-rate
is statistically indistinguishable from chance (95\% CI $[0.30,
0.60]$, straddles 0.5). Iter 43 asks three follow-ups that target
\emph{why} the claim fails:
\begin{enumerate}
  \item \textbf{ZVF-mechanistic decomposition.} Under the i.i.d.
        contrastive-yield model, $ZVF = p^G + (1-p)^G$ should predict
        retention. We compute the residual
        $R_{\rm meas}(G,T) - R_{\rm ZVF}(G,T)$ and test whether the
        $G{=}4$ vs $G{=}32$ gap is concentrated on cells where
        ZVF-implied retention over-predicts the measured value (i.e.,
        the gap is anti-ZVF rather than ZVF-driven).
  \item \textbf{Difficulty stratification.} Among easy prompts (acc
        $\geq 0.75$) and hard prompts (acc $< 0.5$), does the
        $G{=}4$ deficit look the same? Iter 35 hinted that the gap is
        mostly on easy prompts (the contrastive frontier) rather than
        hard ones (signal starvation).
  \item \textbf{Cost-adjusted TOST.} At a fixed FLOP budget, what is
        the smallest equivalence margin $\epsilon$ at which
        $G{=}4 \equiv G{=}32$ is statistically defensible? The
        standard $\epsilon{=}0.02$ may be too tight; larger margins
        ($\epsilon{=}0.10$, $0.20$) reveal whether the gap shrinks
        monotonically or has a plateau.
\end{enumerate}

\paragraph{Effective-equivalence ZVF decomposition.}
We compute the i.i.d. theoretical ZVF at each $(G, T)$ cell using the
cell accuracy as a proxy for $p$. Table~\ref{tab:groupsize-iter43-residual}
shows that \textbf{the measured retention falls below the ZVF-implied
retention for all $G \geq 16$ at every budget} (residuals
$-0.02, -0.13, -0.27$ at $T{=}1$\,M for $G{=}16, 32, 64$ respectively;
$-0.27, -0.34, -0.39$ at $T{=}64$\,M for the same $G$ values). The
sign-flip pattern is clear: for $G \leq 8$, the residual is small
($\leq 0.01$); for $G \geq 16$, the residual grows monotonically with
$G$. This is the central iter 43 finding: \emph{the
$G{=}4$ vs $G{=}32$ retention deficit is not a ZVF-starvation
artifact --- ZVF predicts higher retention than we measure for all
large $G$}.

\begin{table}[htbp]
\centering
\small
\begin{tabular}{lcccc}
\toprule
$(G, T)$ & $R_{\rm meas}$ & $R_{\rm ZVF}$ & Residual & Interpretation \\
\midrule
$G{=}4, T{=}1$\,M  & 0.854 & 0.851 & $+0.004$ & ZVF adequate \\
$G{=}16, T{=}1$\,M & 0.979 & 1.000 & $-0.021$ & ZVF over-predicts \\
$G{=}32, T{=}1$\,M & 0.875 & 1.000 & $-0.125$ & ZVF over-predicts \\
$G{=}64, T{=}1$\,M & 0.729 & 1.000 & $-0.271$ & ZVF over-predicts \\
\midrule
$G{=}4, T{=}64$\,M  & 0.727 & 0.762 & $-0.035$ & modest gap \\
$G{=}16, T{=}64$\,M & 0.965 & 1.000 & $-0.035$ & over-prediction \\
$G{=}32, T{=}64$\,M & 1.000 & 1.000 & $\;\;0.000$ & exact match \\
$G{=}64, T{=}64$\,M & 0.989 & 1.000 & $-0.011$ & near match \\
\bottomrule
\end{tabular}
\caption{\textbf{ZVF-mechanistic residual.}
$R_{\rm meas}$ is measured retention vs the largest-$G$ at the same $T$;
$R_{\rm ZVF}$ is the i.i.d. theoretical retention implied by
$R_{\rm ZVF}(G) = (1 - p^G - (1-p)^G) / \max_G (1 - p^G - (1-p)^G)$.
For $G \geq 16$ the measured retention is consistently below the
ZVF-implied value, indicating that anti-herding alone cannot account
for the $G{=}4$ vs $G{=}32$ gap.}
\label{tab:groupsize-iter43-residual}
\end{table}

\paragraph{Difficulty stratification.}
Binning the $(G, T)$ cells by their accuracy into low ($<0.5$),
mid ($0.5{-}0.75$), and high ($\geq 0.75$) reveals a striking pattern:
\emph{easy prompts (high bin) account for the entire $G{=}4$ deficit}.
At the highest accuracy bin, $G{=}4$ achieves 0.607 (n=3 mid-bin
average), while $G{=}8$ achieves 0.790 (n=2 high-bin), $G{=}16$
achieves 0.835, $G{=}32$ achieves 0.860, and $G{=}64$ achieves
0.835. The gap from $G{=}4$ to $G{=}32$ is $+0.253$ in the
high-accuracy bin and is monotonic in $G$. The mid-bin shows the
same direction but smaller magnitudes ($G{=}4 = 0.607$, $G{=}32 = 0.660$,
gap $+0.053$); the low-bin shows essentially the same ($G{=}4 = 0.410$,
$G{=}32 = 0.420$, gap $+0.010$). This stratification is consistent
with the iter 35 difficulty finding: the $G{=}4$ vs $G{=}32$ gap
\textbf{concentrates on the contrastive frontier where GRPO has the
most signal to lose}.

\paragraph{Cost-adjusted TOST.}
Figure~\ref{fig:group-size-iter43}(c) shows TOST $p$-value curves
across $\epsilon \in [0.005, 0.25]$ for the 10 $(G_a, G_b)$ pairs at
$T{=}64$\,M. The headline results:
\begin{itemize}
  \item At $\epsilon{=}0.02$ (strict equivalence), only \textbf{12/40
        cells} (30\%) are TOST-equivalent --- the Wu claim is too
        tight for the typical $(G_a, G_b, T)$ operating point.
  \item At $\epsilon{=}0.05$, 11/40 cells (27.5\%) pass --- roughly the
        same as the $\epsilon{=}0.02$ cut, indicating the gap is
        larger than 5\% at most operating points.
  \item At $\epsilon{=}0.10$, only 7/40 cells (17.5\%) pass --- the
        claim breaks badly here.
  \item At $\epsilon{=}0.20$, 0/40 cells pass --- no $(G_a, G_b, T)$
        cell is within 20\% accuracy of the large-$G$ benchmark.
\end{itemize}
The Wu 97.6\% retention threshold falls within the 95\% CI at exactly
20/40 cells (50\%, binomial 95\% CI $[0.34, 0.66]$). This is the
sharpest operational verdict: \textbf{half of the (G$_a$, G$_b$, T)
operating space is consistent with Wu's 97.6\% retention; the other
half is not}.

\paragraph{Reproducibility (iter 43).}
The iter 43 artifacts are produced by
\texttt{scripts/group\_size\_iter43.py} (no new measurements: every
number comes from \texttt{group\_size\_token\_normalized.tsv}).
Three TSV outputs are written:
\begin{itemize}
  \item \texttt{group\_size\_iter43\_eff\_zf\_v.tsv} --- 20 rows:
        ZVF-theoretical retention, measured retention, and the
        ZVF-driven residual per $(G, T)$.
  \item \texttt{group\_size\_iter43\_difficulty.tsv} --- 20 rows:
        per-cell difficulty bin (low/mid/high).
  \item \texttt{group\_size\_iter43\_flop\_tost.tsv} --- 40 rows:
        cost-adjusted retention and TOST $p$-values at
        $\epsilon \in \{0.02, 0.05, 0.10, 0.20\}$.
  \item \texttt{group\_size\_iter43\_summary.tsv} --- 13 headline
        metrics for the paper rollup.
\end{itemize}
Plus a 3-panel figure (\texttt{figures/group\_size\_iter43.pdf}).

\paragraph{Frontier synthesis.}
The frontier-model digest (\texttt{FRONTIER\_INSIGHTS.md}) frames ZVF
as ``contrastive yield, not difficulty''; iter 43 is consistent: the
ZVF-mechanistic residual is \emph{negative} for $G \geq 16$, meaning
\emph{contrastive yield over-predicts retention at large $G$}. The
mechanism is consistent with the frontier synthesis's anti-herding
observation: the empirical anti-herding bonus
$\delta_{\rm div} \in [0.13, 0.23]$ (Round 2) is large enough that
ZVF-implied retention reaches 1.0 by $G{=}8$, but the measured
retention only does so by $G{=}32$. The intermediate cells
($G{=}16, 32$ at $T{=}64$\,M) show the residual gap closing as the
real optimisation trajectory catches up to the i.i.d. upper bound.

\begin{figure}[htbp]
\centering
\includegraphics[width=0.95\textwidth]{group_size_iter43.pdf}
\caption{\textbf{Iter 43 Pillar 3 ZVF-decomposed equivalence.}
\textbf{(a)} Measured retention $R(G, T)$ vs Wu 2025's 97.6\% threshold
(red dashed): the threshold is reachable at $T{=}1$\,M for $G{=}8, 16$
but falls below for all $G \geq 32$ at $T \geq 4$\,M.
\textbf{(b)} ZVF-implied retention vs measured retention: the diagonal
identity would mean ZVF alone explains the gap. Points below the
diagonal (large $G$) show the gap is \emph{not} ZVF-driven ---
contrastive yield over-predicts retention.
\textbf{(c)} Cost-adjusted TOST at $T{=}64$\,M: TOST $p$-value vs
equivalence margin $\epsilon$ for all 10 $(G_a, G_b)$ pairs. Only at
$\epsilon \geq 0.20$ does the $p$-value fall below 0.05 for any pair.
Source: \texttt{figures/group\_size\_iter43.pdf}.}
\label{fig:group-size-iter43}
\end{figure}